% 底本 EGAIV.1.pdf，PDF第99--103页，分册印刷第100--104页。
\subsection{域上的代数的情形}
\source{第四卷第一部分，印刷第100--104页；上传PDF第99--103页}
\egas{19.6.1}{定理（Cohen）}
设 $k$ 为域，$K$ 为 $k$ 的扩域，二者均取离散拓扑。则 $K$ 是形式光滑的 $k$-代数，当且仅当扩张 $K/k$ 可分。

必要性将在（19.6.5.1）中证明，在此之前不使用它。这里仅证明充分性，分两种情形。

\noindent\textbf{I．$K/k$ 为有限生成可分扩张。}
由域论（Bourbaki，\textit{代数}，第五章，§9，第3目，定理2），存在纯超越子扩张
$K'=k(T_1,\ldots,T_n)$，使 $K/K'$ 为有限可分代数扩张。由（19.3.5，ii），只须处理纯超越扩张及有限可分代数扩张两种情形。
前一种情形中，$k[T_1,\ldots,T_n]$ 为形式光滑 $k$-代数（19.3.3），故其分式域也形式光滑（19.3.5，iv）。
后一种情形中，对任意与基域作用相容的 $(K,K)$-双模 $L$，均有
\[
 H_k^i(K,L)=0\qquad(i>0).
\]
\cmt{G7-T05}{原文这里说“所有Hochschild上同调群”均为零，须理解为正次数；例如 $K=L=k$ 时 $H_k^0(k,k)=k$。本处证明仅需二次上同调消失，正文写明 $i>0$，不改变后续结论。}
事实上，令 $C=K\otimes_kK$，则 $L$ 是左 $C$-模，而
$H_k^i(K,L)=\Ext_C^i(K,L)$；这里 $K$ 通过乘法同态成为 $C$-模（(M)，IX，4）。因 $K/k$ 有限可分，$C$ 是有限个域的直积，其中通过乘法同态得到的一个直因子就是 $K$（Bourbaki，\textit{代数}，第八章，§8，命题3）。因此 $K$ 是投射 $C$-模，得到上述消失。以 $L$ 为核的每个 $k$-扩张都为 $k$-平凡扩张，尤其交换扩张如此；由（19.4.4）即得结论。

\noindent\textbf{II．一般情形。}
沿用（18.4.5）的记号，只须证明对每个 $K$-向量空间 $L$，有
\[
 \Hom_K\bigl(H_2(P'_\bullet),L\bigr)=0,
\]
亦即 $H_2(P'_\bullet)=0$。若 $K$ 是其子扩张 $K_\alpha$ 的滤过并，则 $P'_\bullet$ 是对应复形 $P'^\alpha_\bullet$ 的归纳极限，因为模的滤过归纳极限与张量积可交换。又因该归纳极限函子正合，
\[
 H_2(P'_\bullet)=\varinjlim_\alpha H_2(P'^\alpha_\bullet).
\]
$K/k$ 可分蕴含每个子扩张可分；而 $K$ 是其有限生成子扩张的滤过并。由第一部分，各右端同调模为零，故 $K$ 形式光滑。\finishproof

\egas{19.6.2}{推论}
设 $A$ 为分离完备的局部环，$K$ 为其剩余域，并\emph{给定一个子域 $k\subset A$}，使通过剩余映射识别后的扩张 $K/k$ 可分。则存在包含 $k$ 的子域 $K'\subset A$，使典范同态 $A\to K$ 在 $K'$ 上的限制为 $K'\xrightarrow{\sim}K$。\cmt{G7-E01}{初版只写“$k$ 为 $K$ 的子域”，未给出相容的 $k\hookrightarrow A$，但结论要求 $K'\subset A$ 包含 $k$，证明也使用 $A$ 的 $k$-代数结构。按初版假设取 $A=\mathbf Z_p$、$K=k=\mathbf F_p$ 即不成立：$A$ 的特征为零，不能包含特征为 $p$ 的子域。正文补明所需嵌入；它在下一推论的应用中本来已经具备。}

设 $\mathfrak m$ 为 $A$ 的极大理想。由（19.6.1），$K$ 为形式光滑的离散 $k$-代数。对完备环 $A$ 及理想 $\mathfrak m$ 应用（19.3.10），恒等同态 $K\to A/\mathfrak m=K$ 提升为 $k$-同态 $f:K\to A$。域到非零环的保单位同态必单射；因此 $f$ 将 $K$ 同构地映到包含 $k$ 的子域 $K'=f(K)$，而剩余映射限制到 $K'$ 是它的逆。

\egas{19.6.3}{推论}
设 $A$ 为完备诺特局部环，$\mathfrak m$ 为其极大理想，$k=A/\mathfrak m$ 为剩余域。以下条件等价：

$a)$ $A$ 包含一个子域；

$b)$ 若 $p=\operatorname{char}k$，则 $pA=0$；

$c)$ 存在域 $K$，使 $A$ 同构于形式幂级数环 $B=K[[T_1,\ldots,T_n]]$ 的一个商环。

这些条件满足时，可以使 $A\simeq B/\mathfrak b$，其中 $\mathfrak b$ 包含于 $B$ 的极大理想的平方。

$c)\Rightarrow a)$：局部环 $A$ 非零，复合同态 $K\to B\to A$ 非零，故单射。
$a)\Rightarrow b)$：若 $K\subset A$ 为子域，其到 $k$ 的剩余同态单射，所以两域特征相同；故 $p\cdot1_A=0$。
$b)\Rightarrow a)$：设 $f:\mathbf Z\to A$ 和 $g:A\to k$ 为典范同态。$\Ker(gf)=p\mathbf Z$，而 $f(p\mathbf Z)=0$，故 $\Ker f=p\mathbf Z$。于是 $A$ 含有同构于 $\mathbf Z/p\mathbf Z$ 的子环 $A'$，且 $A'\cap\mathfrak m=0$。若 $p>0$，$A'$ 已是域；若 $p=0$，则 $A'$ 的每个\emph{非零}元素在 $A$ 中可逆，故 $A$ 还含有它的分式域 $\mathbf Q$。\cmt{G7-T06}{原文称 $A'$ 的“每个元素”可逆，须排除零元。前面的“不与 $\mathfrak m$ 相交”也应按 $A'\cap\mathfrak m=\{0\}$ 理解，不能把零元从两个子环中删除。}

最后证明 $a)\Rightarrow c)$。$A$ 含有与 $k$ 的素子域同构的素子域 $k_0$；$k/k_0$ 可分。由已校订的（19.6.2），存在同构 $f:k\xrightarrow{\sim}K\subset A$，且与剩余映射相容。取 $x_1,\ldots,x_n\in\mathfrak m$，使它们模 $\mathfrak m^2$ 的类为 $\mathfrak m/\mathfrak m^2$ 的 $k$-基。完备性给出唯一连续同态
\[
 u:k[[T_1,\ldots,T_n]]\longrightarrow A,
 \qquad u|_k=f,\quad u(T_i)=x_i.
\]
它由 $x_i$ 的选取而满射（Bourbaki，\textit{交换代数}，第三章，§2，第9目，命题11）。它在常数项和一次相伴分次部分均为同构，故其核包含于 $(T_1,\ldots,T_n)^2$，也得到最后的断言。

\egas{19.6.4}{定理}
设 $k$ 为域，$A$ 为诺特局部 $k$-代数，$\mathfrak m$ 为其极大理想，$K=A/\mathfrak m$，$A$ 取 $\mathfrak m$-预进制拓扑。假设 $K/k$ 可分。则以下条件等价：

$a)$ $A$ 为形式光滑 $k$-代数；

$b)$ $A$ 的完备化 $\widehat A$ 作为拓扑 $k$-代数，同构于 $K[[T_1,\ldots,T_n]]$，其 $k$-代数结构由
$k\to K\to K[[T_1,\ldots,T_n]]$ 给出；

$b')$ $\widehat A$ 作为环同构于某个 $K[[T_1,\ldots,T_n]]$；

$c)$ $A$ 为正则局部环。

$a)\Rightarrow b)$ 是（19.5.4）中 $a)$ 与 $c)$ 等价的特殊情形：将其中的基环换成 $k$，$B$ 换成 $A$，$\mathfrak j$ 换成 $\mathfrak m$，并用（19.6.1）。此时 $\mathfrak m/\mathfrak m^2$ 为有限维 $K$-向量空间，其对称代数是有限变量多项式环。
$b)\Rightarrow b')$ 是直接的；$b')\Rightarrow c)$ 由（17.1.1）得到。
若 $c)$ 成立，由（17.1.1，d），$\gr_{\mathfrak m}^{\bullet}(A)$ 为 $\mathfrak m/\mathfrak m^2$ 的对称代数；再将（19.5.4）中 $b)\Rightarrow a)$ 用于 $C=K$，就得到 $A$ 形式光滑。

\egas{19.6.5}{推论}
设 $k$ 为域，$A$ 为诺特局部 $k$-代数，$\mathfrak m$ 为其极大理想，$A$ 取 $\mathfrak m$-预进制拓扑。若 $A$ 为形式光滑 $k$-代数，则 $A$ 在 $k$ 上几何正则：即对每个有限扩张 $k'/k$，半局部环 $A'=A\otimes_kk'$ 为正则环（17.3.6；另见IV，6.7.6）。此外，若 $K$ 为 $A$ 的剩余域，则 $\widehat A$ 作为环同构于某个 $K[[T_1,\ldots,T_n]]$。

由于 $\widehat{A'}\simeq\widehat A\otimes_kk'$，由（19.3.6）及用于 $A'$ 各极大理想处局部环的（17.1.5），可假设 $A$ 完备。$A'$ 为形式光滑 $k'$-代数（19.3.5，iii），并且是若干完备局部 $k'$-代数的有限直积；每个直因子也形式光滑（19.3.5，v）。所以只须证明一个这样的形式光滑完备局部代数正则。

将其基域仍记为 $k$，令 $k_0$ 为 $k$ 的素子域。$k/k_0$ 可分，故为形式光滑扩张（19.6.1）；再用（19.3.5，ii），$A$ 为形式光滑 $k_0$-代数。其剩余域 $K$ 也在 $k_0$ 上可分，因此对 $A,k_0$ 应用（19.6.4），同时得到正则性和所述形式幂级数环结构。

\egas{19.6.5.1}{推论}
设 $A$ 为诺特局部环，$k\subset A$ 为子域，并假设 $A$ 对其预进制拓扑为形式光滑 $k$-代数。则每个满足 $k\subset K\subset A$ 的子域 $K$ 都是 $k$ 的可分扩张。

事实上，对每个有限扩张 $k'/k$，$K\otimes_kk'$ 单射到 $A'=A\otimes_kk'$ 中。由（19.6.5），$A'$ 正则，因而约化；其子环 $K\otimes_kk'$ 也约化。这就证明 $K/k$ 可分（Bourbaki，\textit{代数}，第八章，§7，第3目，定理1）。\cmt{G7-T07}{原证明末句把由约化性得到的对象写成 $K$；实际须得到 $K\otimes_kk'$ 约化。域 $K$ 本身约化并不能判定扩张可分。正文写出正确对象。}

特别地，对 $A=K$ 应用此推论，便证明了（19.6.1）中的必要性。

\egas{19.6.5.2}{注}
若 $K/k$ 不可分，则配备通常 $k$-代数结构的形式幂级数环 $K[[T_1,\ldots,T_n]]$（其结构同态如（19.6.4））不是形式光滑 $k$-代数。
另一方面，存在形式光滑的完备诺特局部 $k$-代数 $A$，其剩余域 $K$ 在 $k$ 上不可分。例如，取 $B=k[T_1,\ldots,T_n]$，在一个使 $B/\mathfrak n$ 为 $k$ 的有限不可分扩张的极大理想 $\mathfrak n$ 处局部化，然后完备化。这样的 $A$ 不能与通常结构的 $K[[T_1,\ldots,T_n]]$ 为 $k$-代数同构；但由（19.6.5），二者仍可以作为环同构。

某些情形下，可以将（19.6.5）加强如下。

\egas{19.6.6}{命题}
设 $k$ 为域，$A$ 为诺特局部 $k$-代数，$\mathfrak m$ 为其极大理想，并给 $A$ 赋予 $\mathfrak m$-预进制拓扑。假设剩余域 $K=A/\mathfrak m$ 满足：存在有限根式扩张（即有限纯不可分扩张）$k_0/k$，使局部环 $K\otimes_kk_0$ 的剩余域为 $k_0$ 的可分扩张。以后把这一条件称为 $K$ 在 $k$ 上具有有限根式重数（IV，4.7.4）；还将证明，$K/k$ 有限生成时总满足此条件。则以下条件等价：

$a)$ $A$ 为形式光滑 $k$-代数；

$b)$ 存在有限根式扩张 $k'/k$，使 $\widehat A\otimes_kk'$ 作为 $k'$-代数，同构于某个 $K'[[T_1,\ldots,T_n]]$，其中 $K'/k'$ 可分；

$b')$ 存在有限扩张 $k'/k$，使完备半局部环 $\widehat A\otimes_kk'$ 的至少一个局部分量，作为 $k'$-代数同构于这样的 $K'[[T_1,\ldots,T_n]]$，其中 $K'/k'$ 可分；

$c)$ $A$ 在 $k$ 上几何正则（19.6.5）。

先作几点说明。若 $k'/k$ 为有限根式扩张，在 $A'=A\otimes_kk'$ 中只有一个素理想位于 $\mathfrak m$ 之上：它由某个 $p^h$ 次幂属于 $\mathfrak m A'$ 的元素构成，其中 $p$ 为 $k$ 的特征指数（Bourbaki，\textit{交换代数}，第五章，§2，第3目，引理4）。因此 $A'$ 为局部环，$K\otimes_kk'=A'/\mathfrak m A'$ 也为局部环，二者剩余域相同。

若 $K/k$ 可分，则对每个有限扩张 $k''/k$，$K\otimes_kk''$ 为有限个域的直积（Bourbaki，\textit{代数}，第八章，§7，第3目，定理1的推论1）；于是 $\mathfrak m\otimes_kk''$ 是 $A\otimes_kk''$ 的根，而上述各域就是其极大理想处的剩余域。它们都在 $k''$ 上可分，因为对任意扩张 $k'''/k''$，
\[
 (K\otimes_kk'')\otimes_{k''}k'''=K\otimes_kk'''
\]
约化（同处定理1）。将此用于有限根式重数条件，可知对每个有限扩张 $k''/k_0$，$(K\otimes_kk'')_{\mathrm{red}}$ 在 $k''$ 上可分。

现在证明各条件的等价。$b)\Rightarrow b')$ 是直接的。下面证明 $b')\Rightarrow a)$，并明确补足其中的下降步骤。\cmt{G7-G02}{原证明称扩域后的各局部分量与原来的 $\widehat A\otimes_k k'_0$ 同构，这一般错误；例如 $A=k=\mathbf R$、$k'=\mathbf C$ 时，$\mathbf C$ 不与 $\mathbf R$ 为 $\mathbf R$-代数同构。正确结论是各分量彼此由Galois作用联系。正文用这一结论以及（19.4.6）的忠实平坦下降给出修补；初版在此引（19.4.7），与所需判据不符。详细验证见报告G7-G02。}

将给定 $k'$ 放入有限准Galois扩张 $L/k$，并令 $l$ 为 $L$ 中 $k$ 的纯不可分闭包，使 $l/k$ 为有限根式扩张、$L/l$ 为有限Galois扩张。给定的良好局部分量经 $k'\to L$ 扩域后，是有限个形式幂级数环的直积：这是因为 $K'/k'$ 可分，$K'\otimes_{k'}L$ 是在 $L$ 上可分的域的有限直积，且有限扩域与形式幂级数环可交换。因此
\[
 B_0=\widehat A\otimes_kl,\qquad D=B_0\otimes_lL
\]
中的 $D$ 至少有一个这样的局部分量，而 $B_0$ 为完备局部环。

群 $\operatorname{Gal}(L/l)$ 作用在 $D$ 上，并传递地置换 $D$ 的局部分量。确实，$D$ 为完备半局部环，若分量分成不止一个轨道，某个真轨道对应的单位幂等元之和就给出 $D^{\operatorname{Gal}(L/l)}=B_0$ 的非平凡幂等元，矛盾于 $B_0$ 局部。各分量因而作为 $B_0$-代数彼此同构；这些同构对 $L$ 的作用可以是半线性的，但仍保持环的正则性以及剩余域在 $L$ 上的可分性。由（19.6.4），每个分量都为形式光滑的 $L$-代数，故其有限直积 $D$ 也如此（19.3.5，v）。

拓扑 $k$-向量空间 $\widehat A$ 形式投射：一个到离散向量空间的连续线性映射通过某个开子空间商分解，而该商作为 $k$-向量空间投射，因此可连续提升。又因 $L/k$ 为有限自由且忠实平坦的扩张，可应用（19.4.6）的第二种条件，将 $D=\widehat A\otimes_kL$ 的形式光滑性下降到 $\widehat A$。最后由（19.3.6），$A$ 形式光滑。

$a)\Rightarrow c)$ 已由（19.6.5）证明。若 $c)$ 成立，取假设所给的 $k_0$，或包含它的任意有限根式扩张 $k'/k$；则 $A\otimes_kk'$ 为正则局部环，且其剩余域 $K'$ 在 $k'$ 上可分。由（19.6.4），它的完备化
\[
 \widehat{A\otimes_kk'}=\widehat A\otimes_kk'
\]
作为 $k'$-代数同构于 $K'[[T_1,\ldots,T_n]]$，即得 $b)$。
\finishproof

\egas{19.6.7}{注}
（i）在上述证明中，剩余域 $K$ 具有有限根式重数的假设用于 $c)\Rightarrow b)$：它保证能通过有限根式扩域使剩余域可分。修补后的 $b')\Rightarrow a)$ 不需要这一假设。\cmt{G7-T08}{初版此处写成“该假设只用于 $b')\Rightarrow a)$”。这与前一段 $c)\Rightarrow b)$ 中选取 $k_0$ 的实际依赖不符；正文按依赖关系改正，不将该备注误作无条件存在有限根式扩域的保证。}

（ii）以后将在（22.5.8）证明：无须对扩张 $K/k$ 增加条件，$a)$ 与 $c)$ 仍等价。
